\section{ニュートン法}\label{節:ニュートン法}
一次近似は関数の値を求めるだけでなく,方程式の解を探すためにも使える.\,$f(x)=0$の解に近そうな値を$x_n$とし,その点での接線を考えよう.接線は
\[
 y=f(x_n)+f'(x_n)(x-x_n)
\]
である.\,$f'(x_n)\ne0$なら,\,$x$軸との交点は
\[
 x=x_n-\frac{f(x_n)}{f'(x_n)}
\]
となる.この交点を次の値$x_{n+1}$として繰り返す方法を\textbf{ニュートン法}という.曲線上の点から接線を引き,接線の零点から再び曲線へ戻る計算である.

\begin{example}{平方根に近づく数列}
$x_1=2$とし,\,$f(x)=x^2-2$にニュートン法を使って得られる数列を求めよ.また,その数列が$\sqrt2$に収束することを示せ.
\end{example}
$f'(x)=2x$だから,
\[
 x_{n+1}=x_n-\frac{x_n^2-2}{2x_n}
 =\frac12\left(x_n+\frac2{x_n}\right).
\]
順に計算すると,\,$x_2=\frac32$,\,$x_3=\frac{17}{12}$となる.ただし,値が近づいて見えるだけでは収束の証明にはならない.

まず,\,$x_n>\sqrt2$なら,
\[
 x_{n+1}-\sqrt2
 =\frac{x_n^2-2\sqrt2x_n+2}{2x_n}
 =\frac{(x_n-\sqrt2)^2}{2x_n}>0
\]
となる.\,$x_1>\sqrt2$なので,すべての項は$\sqrt2$より大きい.さらに,
\[
 x_{n+1}-x_n=\frac{2-x_n^2}{2x_n}<0
\]
だから単調減少する.下に有界なので極限$L\ge\sqrt2$をもち,漸化式の両辺の極限を取れば
\[
 L=\frac12\left(L+\frac2L\right),\qquad L^2=2
\]
である.\,$L>0$より$L=\sqrt2$と分かる.

途中の式は,誤差の小さくなる様子も表している.\,$e_n=x_n-\sqrt2$とおくと,
\[
 e_{n+1}=\frac{e_n^2}{2x_n}.
\]
例えば$0<e_n<\frac1{10}$なら,\,$x_n>1$なので$e_{n+1}<\frac1{200}$である.接線から作った数列を,差や大小関係を調べる漸化式の問題として考え直せた.

ただし,どの関数でもこのように収束するわけではない.例えば$f(x)=x^3-2x+2$で$x_1=0$とすると,\,$x_2=1$,\,$x_3=0$となり,二つの値を繰り返す.接線の零点を使う理由と,実際に収束することは,それぞれ確かめる必要がある.

